\section{Conclusions and future work}
\label{sec:conclusioni}

\subsection{Conclusions}
\label{sec:conclusioni-sintesi}

This work implemented a two-temperature thermochemical model in OpenFOAM-13
as a basis for simulations of reacting hypersonic flows. The solver advances
the energy equations and uses \mpp{} to recover the temperatures and evaluate
thermochemical properties. The source-term functions are shared with
stand-alone zero-dimensional programs. Where the library's default
formulations differ from those of the reference study, customised source
terms have been implemented using the properties supplied by \mpp.

All zero-dimensional verification cases reported by Casseau
\textit{et al.}~\cite{casseau2016} have been reproduced. Across the cases
computed with both implementations, including reacting cases, the solver
agrees with the zero-dimensional programs to within 0.3\%, except for
high-temperature nitrogen, where the maximum difference is 1.2\%.
Compared with the reference study, final temperatures agree to within 1\%,
nitrogen-mixture composition to within 2.5\% and air composition to within
15\%. The largest temperature discrepancy is a delayed rise in vibrational
temperature in mixtures containing atomic nitrogen. Differences in the
averaging of relaxation times are a plausible explanation, although this
hypothesis remains to be tested.

Reproduction of the results required reconstructing incompletely specified
modelling choices, particularly the treatment of electronic energy and the
reaction-rate constants for air. Verification remains limited to spatially
uniform heat baths with no flow; flow applications require the extensions
outlined below.

\subsection{Future work}
\label{sec:sviluppi}

The proposed extensions address energy transport, viscous-flow modelling and
the remaining discrepancy in vibrational relaxation.

\begin{enumerate}[leftmargin=1.8em]
  \item \textbf{Transport of the vibro-electronic energy.} The vibro-electronic
    energy equation currently contains only the source terms. Convective
    transport must be added, using the same face fluxes already used for the
    species, together with heat conduction and appropriate boundary
    conditions; the source terms should then be evaluated with the updated
    density. The one-dimensional shock tube already included in the project
    provides an initial test case for this extension.
  \item \textbf{Two-temperature viscosity and thermal conductivity.} These are
    required for viscous-flow simulations and wall heat-flux predictions.
    \mpp{} provides these properties and separates thermal conductivity into
    translational and vibro-electronic contributions. The thermophysical
    model must evaluate them in each cell and supply them to the solver in
    place of the one-temperature transport properties.
  \item \textbf{Relaxation-time averaging used in the reference study.}
    Pairwise relaxation times should be evaluated using \mpp{} coefficients
    and averaged as in the reference study. Repeating the molecular-atomic
    nitrogen case would test whether this change accounts for the observed
    delay in vibrational heating.
\end{enumerate}
